\begin{center}
    \Large{\textbf{TÓM TẮT ĐỀ TÀI}}
\end{center}

Hiện nay, việc thiếu hụt các không gian số hỗ trợ tương tác dựa trên vị trí thực tế đang làm giảm đi cơ hội giao lưu và tham gia phong trào của sinh viên. Nhằm tháo gỡ rào cản này, nghiên cứu tiến hành thiết kế và phát triển nền tảng UniConnect. Đây là một giải pháp mạng xã hội tích hợp bản đồ không gian, được định hướng để tối ưu hóa việc khám phá sự kiện và kết nối cộng đồng người học có chung sự quan tâm.

Nền tảng được cấu trúc dựa trên các phân hệ nghiệp vụ trọng tâm: (1) Trực quan hóa không gian số, hỗ trợ người dùng theo dõi và tìm kiếm các hoạt động ngoại khóa, nhóm học tập đang diễn ra tại khu vực lân cận theo thời gian thực; (2) Đề xuất kết nối chuyên sâu, ứng dụng kỹ thuật tìm kiếm vector để đánh giá mức độ tương đồng trong hồ sơ cá nhân, từ đó ghép cặp những sinh viên phù hợp nhất; (3) Quản trị và điều phối sự kiện, cung cấp công cụ để cá nhân tự do khởi tạo và mời những người xung quanh tham gia các hoạt động cục bộ; (4) Trợ lý ảo AI, đóng vai trò như một tư vấn viên hỗ trợ tra cứu thông tin, giải đáp thắc mắc và định hướng lộ trình thông qua giao tiếp bằng ngôn ngữ tự nhiên; (5) Phân quyền và bảo vệ dữ liệu, cho phép thiết lập linh hoạt các bán kính hiển thị và ẩn danh vị trí để đảm bảo an toàn thông tin.

Về mặt nền tảng công nghệ, dự án tập trung hoàn thiện kiến trúc máy chủ (Backend) với khả năng xử lý truy vấn dữ liệu đa chiều hiệu năng cao. Việc đóng gói hệ thống qua môi trường ảo hóa cũng được thực hiện nhằm đảm bảo tính sẵn sàng cho các pha triển khai tiếp theo. Kết quả của đồ án không chỉ cung cấp một mô hình kỹ thuật phần mềm vững chắc mà còn mang lại giá trị thực tiễn cao, góp phần kiến tạo một môi trường đại học năng động và gắn kết hơn.